% Exact finite-input counterexample: both marginal inputs Bernoulli(1/2).
\begin{tikzpicture}[foundation picture]
 \path[use as bounding box] (0,0) rectangle (160,83);
 \draw[FigureBlue,rounded corners=1.8pt] (1,25) rectangle (38,68);
 \node at (19.5,64) {可观察的源数据};
 \node at (19.5,56) {$x=0\longmapsto y=0$};
 \node at (19.5,48) {$x=1\longmapsto y=1$};
 \draw[FigureMuted,dashed] (2,42)--(37,42);
 \node at (19.5,37) {无标签目标输入};
 \node at (19.5,30) {$x=0,\ 1$};
 \node[foundation node,draw=LancetGreen,minimum width=23mm,minimum height=12mm] (alg) at (57,47) {算法 $A$};
 \node[foundation node,draw=LancetGreen,minimum width=24mm,minimum height=12mm,align=center] (hyp) at (86,47) {同一输出\\$h(x)=x$};
 \draw[foundation flow] (38,47)--(alg.west);
 \draw[foundation flow] (alg.east)--(hyp.west);
 \node[foundation note] at (57,36) {固定内部随机数};
 \foreach \y/\title/\casecolor/\ya/\yb/\risk in {50/{y=x}/LancetGreen/0/1/0,5/{y=1-x}/FigureRed/1/0/1}{
  \draw[\casecolor,rounded corners=1.8pt] (107,\y+4) rectangle (159,\y+33);
  \draw[\casecolor] (107,\y+25)--(159,\y+25);
  \node at (133,\y+29) {隐藏目标规律 $\title$};
  \node at (116,\y+20) {输入};\node at (132,\y+20) {预测};\node at (149,\y+20) {真值};
  \node at (116,\y+14) {$0$};\node at (132,\y+14) {$0$};\node at (149,\y+14) {$\ya$};
  \node at (116,\y+8) {$1$};\node at (132,\y+8) {$1$};\node at (149,\y+8) {$\yb$};
  \node at (133,\y-1) {$R_T(h)=\risk$};
 }
 \draw[foundation flow] (hyp.east)--(102,47)|-(107,65);
 \draw[foundation flow] (hyp.east)--(102,47)|-(107,20);
\end{tikzpicture}
